% !TeX program = tectonic
\documentclass[11pt,a4paper]{article}
\usepackage{fontspec}
\usepackage[english]{babel}
\babelfont{rm}[Language=Default]{Liberation Serif}
\babelfont[Language=Default]{sf}{Carlito}
\babelfont[Language=Default]{tt}{DejaVu Sans Mono}
\usepackage{amsmath,amssymb,amsthm}
\usepackage{geometry}
\geometry{a4paper, top=2.4cm, bottom=2.4cm, left=2.4cm, right=2.4cm}
\usepackage{booktabs,tabularx,enumitem,xcolor,tcolorbox}
\tcbuselibrary{breakable,skins}
\definecolor{linkblue}{HTML}{1F4E79}
\definecolor{statgreen}{HTML}{2F6B3A}
\definecolor{goldacc}{HTML}{8B7E5A}
\usepackage[colorlinks=true,linkcolor=linkblue,urlcolor=linkblue,unicode=true]{hyperref}
\theoremstyle{plain}
\newtheorem{theorem}{Theorem}
\newtheorem{lemma}[theorem]{Lemma}
\newtheorem{proposition}[theorem]{Proposition}
\newtheorem{corollary}[theorem]{Corollary}
\theoremstyle{definition}
\newtheorem{definition}[theorem]{Definition}
\theoremstyle{remark}
\newtheorem{remark}[theorem]{Remark}
\newtcolorbox{statusbox}[1][]{colback=statgreen!5,colframe=statgreen!75,
  title={\textbf{Summary of results:} #1},fonttitle=\small,boxrule=0.7pt,arc=2pt,breakable}
\newtcolorbox{databox}[1][]{colback=goldacc!6,colframe=goldacc!80,
  title={\textbf{Protocol data:} #1},fonttitle=\small,boxrule=0.7pt,arc=2pt,breakable}
\newtcolorbox{verifybox}[1][]{colback=linkblue!5,colframe=linkblue!80,
  title={\textbf{Verification:} #1},fonttitle=\small,boxrule=0.7pt,arc=2pt,breakable}
\widowpenalty=10000
\clubpenalty=10000
\setlength{\parskip}{0.35em}
\newcommand{\CC}{\mathbb{C}}
\newcommand{\RR}{\mathbb{R}}
\newcommand{\ZZ}{\mathbb{Z}}
\newcommand{\QQ}{\mathbb{Q}}
\newcommand{\Ga}{\Gamma}
\newcommand{\Bfun}{\mathrm{B}}
\newcommand{\im}{\operatorname{Im}}
\newcommand{\re}{\operatorname{Re}}
\newcommand{\lcm}{\operatorname{lcm}}
\newcommand{\SNF}{\operatorname{SNF}}
\newcommand{\Jac}{\operatorname{Jac}}
\newcommand{\rank}{\operatorname{rank}}
\newcommand{\Ind}{\operatorname{Index}}
\newcommand{\disc}{\operatorname{disc}}
\begin{document}
\begin{center}
{\LARGE\bfseries Certificate F: Rationalization}\\[0.4em]
{\large the a priori denominator bound and exact LLL}\\[0.6em]
{\normalsize Isaev Iskhak Khamzatovich}\\[0.2em]
{\normalsize The <<Dynamic Principle>> Program --- the \texttt{hodge-laboratory}}\\[0.15em]
{\small Standalone theorem monograph · Монография 14/16}
\end{center}
\vspace{0.4em}
{\color{linkblue}\hrule height 0.8pt}
\vspace{0.8em}

\section{Setting}

Certificate F solves rationalization: recover exact rational values
from numerical measurements with an a priori uniqueness guarantee.
The key is the denominator bound computed before the measurements.

\begin{theorem}[F1--F4]\label{th:Fstand}
\begin{enumerate}[leftmargin=2em]
\item[(F1)] Uniqueness: if $2\varepsilon<1/Q^2$ with $Q$ the
denominator bound, an interval of width $2\varepsilon$ contains at
most one rational with denominator $\le Q$. At $100$ digits the
certified bound is $Q<7.07\cdot10^{49}$.
\item[(F2)] The a priori bound is computed from geometry: K3 ---
$4|\det G_8|=256$ (Cramer), torus --- $1$, Klein --- $2$ (in the
$\tau$-basis) and $1$ (over integers).
\item[(F3)] The integer $q$-scan recovers values without floating
point: complexity $O(nQ)$, acceptance $d_q\le q/2$.
\item[(F4)] Exact LLL: $\im\tau=\sqrt7/2$ recovered with the minimal
polynomial $4X^2-7$; degree-2 polynomials are unique.
\end{enumerate}
\end{theorem}

\begin{proof}
(F1) The diophantine estimate: two distinct reduced fractions with
denominators $\le Q$ differ by at least $\frac1{qQ'}\ge\frac1{Q^2}$;
an interval of width $<1/Q^2$ cannot contain both.
(F2) By Cramer's rule the target coordinates are ratios of
determinants built from the exact intersection matrix; denominators
divide $\det G_8=64$; the bound $4\cdot64=256$ has margin.
(F3) For each candidate $q\le Q$: $d_q=\min_n|qx-n|$ in integer
multiplications; accept at $2d_q\le1$.
(F4) LLL on the Klein period lattice: exact operations in
$\QQ(\sqrt{-7})$ give $\im\tau=\sqrt7/2$; the polynomial $4X^2-7$ is
unique; $j(\tau)=-3375$ with deviation $1.99\cdot10^{-117}$.
\end{proof}

\begin{databox}[the end-to-end K3 run]
$49$ measurements $\to$ recovery $[Z]=[target]$ exact; maximal
denominator $4$; the a priori bound $Q=256$ (margin $\times64$);
torus: degrees $(2,2)$; Klein: $\re\tau=1/2$, $\im\tau=\sqrt7/2$,
$\mathrm{tr}\,t=-1$, $\mathrm{norm}\,t=2$, $j=-3375$; the
$\im\tau$ residual $7.75\cdot10^{-121}$. Negative test: the value
$4\pi^2$ is \emph{rejected} by the scan --- consistent with
Lindemann transcendence.
\end{databox}

\begin{remark}[the role of the negative test]
The rejection of $4\pi^2$ is a power control: the scan does not
accept ``similar'' values. The agreement with Lindemann's theorem
(the transcendence of $\pi$) independently confirms the correctness
of the acceptance criterion.
\end{remark}

\begin{statusbox}[summary]
Certificate F accepted ($5/5$): rationalization with a guaranteed
bound works end-to-end; the negative test passes.
\end{statusbox}

\begin{verifybox}[reproduction]
\texttt{python3 laboratory.py} $\to$ ``Certificates $\to$ F''; the
designer: an arbitrary value and bound $Q$ --- the lab recovers and
prints the verdict.
\end{verifybox}

\vfill
{\color{linkblue}\hrule height 0.6pt}
\vspace{0.5em}
{\small\color{gray}
License: individual exclusive license of Isaev Iskhak Khamzatovich.
All rights to the computations, formulas, and texts belong to the author.
Verification: \texttt{python3 laboratory.py} (``Stands''/``Certificates''),
Lean: \texttt{verification/lean/HodgeLaboratory.lean}.
}
\end{document}
